\section{Lo que enseña el catálogo}
\label{sec:limites:catalogo}

Las Tablas~\ref{tab:limites-rates} y~\ref{tab:limites-rates-b} recogen el caudal, la cuota y la ventana de cada API, y de ellas salen seis afirmaciones sobre la relación entre límites de uso y \textit{rates}.
Cada afirmación se cuenta sobre las ocho APIs y no sobre sus planes, para que una API con cinco planes no pese cinco veces.
El recuento reparte las ocho en cuatro grupos, que se escriben siempre en este orden: a favor, en contra, no comprobable y no aplica.
Una API cuenta a favor cuando su Datasheet tiene el dato que la afirmación necesita y ese dato va en la dirección enunciada.
Cuenta en contra cuando tiene el dato y apunta al lado opuesto, y entonces es un contraejemplo y se nombra.

Los otros dos grupos separan el dato que falta de la afirmación que no se puede plantear.
Una API es no comprobable cuando la afirmación tiene sentido en ella pero ningún hecho recogido aporta el dato, y en ese caso nunca suma a favor.
No aplica cuando la afirmación presupone una estructura que la API no tiene, como los planes en Mailchimp Transactional o la cuota incluida en Amazon SES.

El cociente $R \cdot W / Q$ es el de la Sección~\ref{sec:limites:geometrica}: cuántas veces cabe la cuota en lo que el \textit{rate} admite durante su ventana.
Como se calcula con un destinatario por petición ($k = 1$), es una cota inferior, porque enviar a varios destinatarios en cada petición solo puede agrandarlo.

La primera afirmación es que, con la cuenta en estado sano, la cuota es la restricción activa en toda API cuya Datasheet lleva \textit{rate} y cuota con valor.
El recuento es de seis a favor, ninguna en contra, una no comprobable y una que no aplica.
El cociente $R \cdot W / Q$ va de 103,7 en MailerSend Starter a 129.600.000 en el plan Free de Postmark, es decir, de dos a ocho órdenes de magnitud.
Con los volúmenes contratados que el proveedor publica y la Datasheet no recoge, el mínimo baja a 2,59 en MailerSend Professional 2M, algo más del doble.

La no comprobable es Amazon SES, porque fuera del entorno de pruebas el proveedor no publica ni el \textit{rate} ni la cuota, que cada cuenta recibe al obtener el acceso de producción.
La que no aplica es Mailchimp Transactional, que no tiene planes y cuya franquicia de prueba de 500 correos no caduca ni se renueva.
Sin ventana en la cuota, el cociente no está definido.
Con 25 correos por hora, agotar esa franquicia lleva veinte horas, así que en esta API el tiempo lo pone el caudal.

En esas seis APIs el \textit{rate} del cociente sale de fuentes distintas, recogidas en la Sección~\ref{sec:limites:publicado}.
Solo MailerSend lo publica para todos sus planes, con 120 peticiones por minuto, y Resend publica un valor por defecto por equipo.
Postmark lo devuelve en la cabecera de respuesta, y la de Mailgun da 80.000 peticiones por una ventana medida de unos diez segundos, que la Tabla~\ref{tab:limites-rates-b} divide para obtener 8.000 por segundo.
SendGrid lo da en unas preguntas frecuentes sin plan, y Brevo no dice qué plan cae en cada categoría.
Para dimensionar un volumen de envío, el \textit{rate} puede ignorarse en seis de las ocho APIs mientras la cuenta esté sana, y en Amazon SES la pregunta no tiene respuesta antes de abrir la cuenta.

La segunda afirmación tiene dos mitades que se cuentan aparte: el \textit{rate} del envío individual no varía con el plan, y lo que el plan compra es cuota.
La primera mitad se limita al envío individual porque MailerSend publica para el envío por lotes, \texttt{/v1/bulk-email}, un caudal que sube con el plan: 10 peticiones por minuto en Hobby, 15 en Starter y 30 en Professional.
Esa mitad queda en tres a favor, una en contra, dos no comprobables y dos que no aplican.
Las tres a favor son Resend, SendGrid y MailerSend.
Resend publica un valor por defecto por equipo, SendGrid publica 10.000 peticiones por segundo sin condición de plan y MailerSend publica 120 peticiones por minuto en \texttt{/v1/email} para sus cuatro planes.

El contraejemplo es Brevo, con las tres categorías de caudal para el mismo \textit{endpoint} de envío que recoge la Sección~\ref{sec:limites:publicado}.
Pasar de la categoría General a la Extended multiplica por seis el caudal del envío individual, que así puede crecer al subir de plan.
El dato va en ámbar (\semambar{}) porque la página no dice qué plan cae en cada categoría.
Postmark y Mailgun son no comprobables, porque su caudal de envío se observó en la cabecera de una sola cuenta gratuita y la Datasheet lo extiende a todos los planes por falta de dato.
La mitad no aplica a Amazon SES, cuyo único caudal publicado es el del entorno de pruebas, un estado de la cuenta y no un plan.
Tampoco aplica a Mailchimp Transactional, que no tiene planes y fija el caudal según su reputación.

La segunda mitad queda en seis a favor y dos que no aplican.
Las seis a favor son Postmark, Resend, SendGrid, MailerSend, Brevo y Mailgun, cuya cuota crece con el precio del plan, en Brevo de 300 correos al día a diez millones al mes.
En Brevo el plan compra además caudal de envío individual, y en MailerSend, caudal de envío por lotes.
En Amazon SES el plan no compra volumen, porque ningún plan incluye cuota: los 200 correos al día del entorno de pruebas valen igual para los cuatro planes, que compran funciones y tramos de precio unitario \citep{sespricing}.
En Mailchimp Transactional no hay planes, y lo que se compra es la propia unidad de volumen, el bloque de 25.000 correos, a un precio que baja de 20 a 10 USD por bloque \citep{mailchimppricing}.

La tercera afirmación es que el \textit{rate} que publica el proveedor coincide con el que devuelve la API en sus cabeceras o con el que mide un sondeo.
El recuento es de cuatro a favor, ninguna en contra, cuatro no comprobables y ninguna que no aplique.
Las cuatro que coinciden publican el valor con su ámbito escrito: Resend lo da por equipo, MailerSend por estado de la cuenta, Brevo ruta a ruta y Mailchimp Transactional para la cuenta de prueba.
La cabecera del envío de Brevo devolvió las 1.000 peticiones por segundo de su tabla General, y en Mailchimp Transactional la documentación, el panel y la medición coinciden en 25 correos por hora.

En MailerSend la coincidencia es de valor y no de estado.
Las 10 peticiones por minuto que el proveedor publica para la cuenta en revisión son las que el sondeo midió en una cuenta sin aprobar, que es otro estado, y las 120 de la cuenta sana siguen sin observar.

Dos de esas coincidencias tienen una segunda cifra al lado.
En Brevo la coincidencia vale solo para la categoría General, porque su guía de envío por lotes da para el mismo \textit{endpoint} 6.000 peticiones por hora, 600 veces menos.
En Resend, la Datasheet de ejemplo que publica el repositorio de la guía de ACM2YAML, fechada en febrero de 2026, fija 2 peticiones por segundo \citep{apisconsumptionmodel}, y la documentación consultada en agosto de 2026 publica 10 \citep{resenddocs}.
Las capturas de su página de límites en Wayback Machine muestran que cambió la oferta: 2 por segundo hasta marzo de 2026, 5 entre abril y julio, y 10 desde agosto.

Las cuatro no comprobables carecen de una cifra publicada y observada para el mismo \textit{endpoint} y el mismo ámbito.
Postmark no publica \textit{rate} y devuelve 5.000 en la cabecera de cada respuesta.
SendGrid publica 10.000 peticiones por segundo para un \textit{endpoint} de envío que, en modo \textit{sandbox}, no devuelve cabeceras, y en los de lectura devuelve 600 y 6 en \texttt{X-RateLimit-Limit}, valores que ninguna página documenta.
Amazon SES no publica su caudal de producción.

En Mailgun, la cifra publicada y la medida tienen ámbitos distintos.
Sus preguntas frecuentes dan 1.500 peticiones por segundo por IP de cliente para toda la API, y la cabecera del envío cuenta otro contador, 80.000 peticiones por una ventana de unos diez segundos (unas 8.000 por segundo si se dividen).
La medición no dice cuál de los dos rige en el envío, así que no contradice la cifra publicada.

La cabecera de respuesta tampoco es una señal fiable en todo el catálogo.
Amazon SES no emite ninguna cabecera de caudal, y sus nueve respuestas 429 medidas no traen \texttt{Retry-After}, la cabecera que dice cuánto esperar antes de reintentar.
Mailchimp Transactional no emite cabecera de caudal en ninguna respuesta del sondeo ni la menciona en ninguna página oficial.
Mailgun la emite solo en un POST aceptado, así que leer el caudal le cuesta al cliente un envío.

Donde la cabecera llega, la documentación no siempre la describe bien.
Brevo afirma que llega en todas las respuestas, y en la medición llegó en una respuesta de error 400 y no en dos respuestas correctas 200.
En SendGrid, \texttt{X-RateLimit-Reset} trae en tres \textit{endpoints} un valor de una o dos cifras (6, 9 y 58), que se lee como los segundos que faltan, y en otros tres una marca de tiempo Unix.
La documentación solo muestra la segunda forma \citep{sendgridratelimits}.
Una herramienta que lea esa cabecera para calcular cuándo reintentar encuentra el mismo nombre con dos formas de valor y ningún aviso.

La alternativa a la cabecera es un \textit{endpoint} que informa de la capacidad restante de la cuenta sin gastar cuota, y cuatro de las ocho APIs lo tienen, aunque el de Brevo da el saldo y no el límite.
Amazon SES devuelve en un \texttt{GET} la cuota de 24 horas y el caudal máximo de envío de la cuenta, y su documentación lo señala como la vía para conocerlas.
Mailchimp Transactional devuelve la cuota horaria junto a la reputación en \texttt{POST /users/info}, y MailerSend tiene \texttt{GET /v1/api-quota} \citep{mailersenddocs}.
El \texttt{GET /v3/account} de Brevo devuelve el saldo de correos y el periodo exacto del ciclo, pero ninguno de sus campos es de límite.

Postmark, Resend, SendGrid y Mailgun no tienen ninguno, y en ellas el cliente conoce su capacidad restante después de gastarla o no llega a conocerla.
Donde existe, el \textit{endpoint} tampoco garantiza un dato al día.
En MailerSend, once horas y veinticuatro minutos después de la medianoche que repone la cuota diaria, seguía sirviendo el saldo del día anterior.

La cuarta afirmación es que el proveedor publica un tope de destinatarios por petición, el valor máximo del factor $k$ con el que se pasa del \textit{rate}, que cuenta peticiones, a la cuota, que cuenta correos.
El recuento es de siete a favor, una en contra, ninguna no comprobable y ninguna que no aplique, porque siete APIs publican ese tope y una no.
Postmark, Resend y Amazon SES admiten 50 destinatarios por petición, MailerSend 70, SendGrid 1.000, Mailgun 1.000 por lote y Brevo 2.000.
El contraejemplo es Mailchimp Transactional, que escribe que no hay limitación de destinatarios y que el JSON de la llamada debe pesar menos de 10 MB \citep{mailchimpdocs}.
Ahí la cota de una llamada está en megabytes, y no hay $k$ que pase de una unidad a la otra.

Cinco de los siete topes están verificados por sondeo: los de Postmark, Resend, SendGrid y Mailgun, y el de la cuenta sin aprobar de MailerSend.
En Mailgun el destinatario 1.001 devuelve un error 400 que nombra el máximo de 1.000.
En Brevo no hubo sondeo de destinatarios, y el tope de Amazon SES sale de su página de cuotas, que lo define como la suma de \texttt{To}, \texttt{Cc} y \texttt{Bcc} \citep{sesdocs}.
En Amazon SES una petición de 50 destinatarios consume la cuarta parte de la cuota diaria del entorno de pruebas, y en Brevo una de 2.000 equivale a casi siete veces la cuota diaria del plan Free.

El factor $k$ solo conecta las dos magnitudes si cada destinatario descuenta una unidad de cuota.
De las cuatro APIs en las que esa equivalencia está sostenida (Sección~\ref{sec:limites:geometrica}), Resend y SendGrid la publican, y en Brevo y Amazon SES salió de medir.
En Brevo, el saldo de la cuenta bajó de 300 a 298 tras una petición con dos destinatarios.
En Amazon SES, el contador marcó 5 sobre 200 tras cuatro mensajes con cinco destinatarios, porque uno llevaba la misma dirección en \texttt{to} y en \texttt{cc}.
Ese recuento descarta a la vez que la unidad sea el mensaje y que el proveedor deduplique direcciones, porque cualquiera de las dos habría dado cuatro.
En Postmark, MailerSend, Mailgun y Mailchimp Transactional esa equivalencia es una suposición que ningún dato confirma.

La quinta afirmación es que el proveedor no publica un caudal propio de cada plan ni cuándo se repone la cuota.
En la mitad que trata del caudal del envío individual, el recuento es de siete a favor, una en contra, ninguna no comprobable y ninguna que no aplique.
Un caudal común a todos los planes, como las 120 peticiones por minuto de MailerSend, cuenta a favor, porque no es propio de ningún plan.
El contraejemplo es Brevo, la única que pone una cifra de ese caudal junto a una categoría de plan.
Amazon SES cuenta a favor porque no publica el caudal de producción.
El envío por lotes de MailerSend, con 10, 15 y 30 peticiones por minuto según el plan, queda fuera de esta mitad por el mismo motivo que en la segunda afirmación.

Cuatro de las siete Datasheets declaran de formas distintas el caudal que el proveedor no liga a un plan.
La de Amazon SES deja vacío el bloque de caudal, la de Mailgun referencia un único valor desde sus cuatro planes y la de MailerSend toma las 120 peticiones por minuto como valor por defecto.
La de Mailchimp Transactional repite el caudal como estado de la cuenta en sus tres \textit{endpoints} de envío.

La otra mitad de la quinta afirmación, la de la ventana de la cuota, queda en ocho a favor.
Ninguna API tiene esa ventana en verde si cada una se clasifica por el peor semáforo de sus cuotas: cuatro están en rojo (\semrojo{}) y cuatro en ámbar (\semambar{}).
SendGrid es la que más se acerca, porque publica la medianoche UTC y el periodo de facturación, y solo la ventana de su plan Premier queda en ámbar.

Los cuatro rojos responden a tres causas.
Mailgun y Brevo dan valor y periodo para las cuotas de su plan Free, pero ninguna fuente fecha su reposición.
Amazon SES describe su cuota diaria como deslizante en una página, es decir, que en cada envío cuenta lo enviado en las 24 horas anteriores, y como fija en otra, en la que hay que esperar al siguiente periodo de 24 horas.
El campo se deja vacío a propósito.
Mailchimp Transactional escribe que su franquicia no caduca ni se renueva, mientras ACM2YAML v0.3 exige un periodo en toda cuota.
En esas cuatro APIs, quien integra la API no puede calcular cuándo se repone su cuota sin observarlo en su propia cuenta.

La ventana del caudal no la declara ninguna de las ocho, y dos sondeos la acotaron.
En Brevo la ruta \texttt{GET /v3/smtp/emails}, que consulta la actividad de envío, volvió a admitir peticiones antes de cumplirse un segundo desde la primera de la serie, y con una sola muestra no se distingue una ventana fija alineada con el reloj de una reposición gradual.
En Mailgun la ventana dura entre 9,85 y 9,87 segundos, la abre la primera petición y no es deslizante, y ninguno de los tres tipos de ventana que admite ACM2YAML v0.3 (deslizante, fija a la medianoche UTC y fija al ciclo de facturación) la describe.

La sexta afirmación es que el estado de la cuenta condiciona la capacidad, con independencia del plan.
El recuento es de seis a favor, dos en contra, ninguna no comprobable y ninguna que no aplique.
Son las seis APIs en que la capacidad depende del estado de la cuenta y no del plan, recogidas como F-18 en el Apartado~\ref{subsec:evolucion:p06}.
Ese apartado describe el caso de MailerSend y la extensión que se propone para modelar el estado de la suscripción.
Solo MailerSend da una cifra para cada estado: 120 peticiones por minuto en \texttt{/v1/email} con la cuenta sana y 10 con la cuenta en revisión.

Las otras cinco a favor son Amazon SES, Mailchimp Transactional, Postmark, Resend y Mailgun.
Amazon SES publica los topes del entorno de pruebas y no los de producción.
Mailchimp Transactional fija 25 correos por hora en la cuenta de prueba y ata la de pago a su reputación.
Postmark limita la cuenta en modo de prueba a 100 correos y a dominios verificados, y con la cuenta pendiente de aprobación rechazó con el código 412 cada mensaje a otro dominio.
Resend sube su valor por defecto de 10 peticiones por segundo a los remitentes de confianza, sin decir cuánto ni qué umbral los define.
Mailgun eleva el caudal a medida que crece la reputación de envío, sin una sola cifra.
Ninguna de las seis publica el umbral que dispara el cambio de estado.

Las dos en contra son SendGrid y Brevo.
En SendGrid el único estado que cambia la capacidad es el fin de la prueba gratuita, que es la caducidad de un plan y no un estado de la cuenta.
En Brevo lo que cambia el caudal es la categoría del acuerdo comercial y no el estado de la cuenta, y para la cuenta sin aprobar no hay cifra.
ACM2YAML v0.3 ya cubre parte de esos estados con sus alias de estado, nombres bajo un \textit{endpoint} que agrupan los límites de cada estado, como el \texttt{sandbox} de Amazon SES (Apartado~\ref{subsec:evolucion:p06}).

Fuera de las seis afirmaciones, el catálogo muestra qué pasa con el envío que supera la cuota.
Cinco de las ocho APIs cobran un sobrecoste por cada correo que pasa de la cuota del plan, y en las cinco su precio baja con el volumen.
Los tramos de Postmark y Resend ya se recogieron en el Apartado~\ref{subsec:evolucion:p01}, donde un mismo plan cambia de sobrecoste con el volumen contratado.
En las otras tres, MailerSend va de 1,50 a 0,50 USD por cada mil correos y Mailgun de 1,80 a 0,40 USD \citep{mailgunpricing}.
SendGrid fija el sobrecoste por correo, y llevado a mil correos va de 1,30 a 0,50 USD.

Las tres APIs sin sobrecoste ponen algo distinto en su lugar.
Brevo niega que haya cargos adicionales y sube de plan automáticamente cuando crece el número de contactos \citep{brevopricing}.
Amazon SES no tiene cuota incluida que exceder, y sus tramos de precio bajan con el volumen sin corte ni penalización \citep{sespricing}.
Mailchimp Transactional detiene el envío al agotar el bloque, y lo que ocupa el lugar del sobrecoste es comprar otro.

En el borde de la cuota hay tres comportamientos documentados.
Amazon SES rechaza la llamada que pasa de la cuota de 24 horas de la cuenta y descarta el mensaje sin ponerlo en cola ni reintentarlo.
SendGrid en su prueba gratuita, MailerSend en el plan Free y Mailchimp Transactional al agotar los bloques detienen el envío.
Brevo retiene en cola hasta 1.000 correos por encima de la cuota diaria y no entrega los que pasan de esa cifra.
Los Apartados~\ref{subsec:evolucion:p07} y~\ref{subsec:evolucion:p09} recogen las dos propuestas que salen de este borde: el tope de diez veces el volumen del plan que SendGrid impone al sobrecoste y la cola con la que Brevo retiene los correos tras agotar la cuota.

Ninguna de las seis afirmaciones se cumple entera en las ocho APIs, y Mailchimp Transactional y Amazon SES quedan fuera de al menos dos de ellas.
